\subsection{Basic notation}
\label{sec:zinc-crypto-preliminaries}

\albert{Check how much of this needs to be given here } 
%%%%%%%%%%%%%%%%%%%%%%%%%%%%%%%%%%%%%%%%%%%%%%%%%%%%%%%%%%%%%%%%%%%%%%%%%%%%%%%%
%%%%%%%%%%%%%%%%%%%%%%%%%%%%%%%%%%%%%%%%%%%%%%%%%%%%%%%%%%%%%%%%%%%%%%%%%%%%%%%%


For $k\geq 1$, let $[k]\defeq \{1,\ldots, k\}$, and let $\range{0}{k}$ be $\{0\}\cup[k]$. We write $x\getsr S$ for uniform sampling from a finite set $S$, and $x\gets\mathcal{D}$ for sampling from a distribution $\mathcal{D}$.


\albert{we are not doing that?} Vectors and tuples are denoted by boldface letters $\vv, \uu, \ldots$, with $\vv_i$ the $i$-th entry. When $\vv$ denotes a matrix, $\vv_{i}$ is the $i$-th row and $\vv_{i,j}$ the $(i,j)$-th entry.
%For execution traces (cf.\ \cref{s:uair}), we write $\vv_{\colindex}[\rowindex]$ for the entry at row $\rowindex$ and column $\colindex$, and $\vv[\rowindex]$ for the row vector at row $\rowindex$.
Given a vector $\vv$ of length $n$, we write $\set(\vv)\defeq\{\vv_i:i\in[n]\}$ for the set of entries of~$\vv$. %For an $m\times n$ matrix $M$ and a column index $j\in[n]$, we write $\set(M_{\cdot,j})\defeq\{M_{i,j}:i\in[m]\}$ for the set of entries in column~$j$.



%%%%%%%%%%%%%%%%%%%%%%%%%%%%%%%%%%%%%%%%%%%%%%%%%%%%%%%%%%%%%%%%%%%%%%%%%%%%%%%%
%%%%%%%%%%%%%%%%%%%%%%%%%%%%%%%%%%%%%%%%%%%%%%%%%%%%%%%%%%%%%%%%%%%%%%%%%%%%%%%%

We use $\KK$ to denote (possibly infinite) fields and $R$ to denote rings. We write $\FF_q$ for the finite field with $q$ elements (where $q$ is a prime power), $\QQ$ for the field of rational numbers, and $\ZZ$ for the ring of integers. We use standard terminology and notation: an ideal $I$ generated by a set $S\subseteq R$ is denoted $(S)$, the quotient ring is denoted $R/I$. A ring is an integral domain if the product of any two nonzero elements is nonzero. 



\subparagraph{Polynomial rings.}
For variables $\YY=(Y_1,\ldots, Y_\mu)$, $R[\YY]$ denotes multivariate polynomials over $R$, and $R^{\multilin}[\YY]$ the multilinear polynomials (individual degree $\le 1$ in each variable).

We often work with multivariate polynomials whose coefficients are themselves univariate polynomials. For a ring $\cR$ (typically $\ZZ[X]$, $\QQ[X]$, or $\local{q}[X]$) and variables $\YY=(Y_1,\ldots, Y_\mu)$, we write $\cR[\YY]$ for polynomials in $\YY$ with coefficients in $\cR$. We view elements of $(\ZZ[X])[\YY]$ as polynomials in $\YY$ (not as elements of $\ZZ[X,\YY]$), denoting them by $f(\YY)$.
For polynomial rings, $\cR^{<d, <B}[X]$ denotes polynomials of degree $<d$ with coefficients in $\cR^{<B}$; we may write $\cR^{<d}[X]$ when the bit bound is clear.
Ring homomorphisms $f:\cR_1 \to \cR_2$ extend coefficient-wise to $\cR_1[\YY] \to \cR_2[\YY]$.
More generally, all maps are extended entrywise to vectors and matrices.


% write $I \cdot \FF_{q^\ell}[X]$ for its \emph{scalar extension} to $\FF_{q^\ell}[X]$, i.e., the ideal generated by $I$ in $\FF_{q^\ell}[X]$.

\subparagraph{Localization of rationals at a prime.}
\label{sec:zinc-crypto-local-rings}
The subset $\loc{p} = \{a/b\in \QQ : p\nmid b\}$ of $\QQ$ is a ring called the \emph{localization of $\ZZ$ at the prime~$p$}. This ring admits a surjective homomorphism $\phi_p: \loc{p} \to \FF_p$, $a/b \mapsto a \cdot b^{-1}$, with kernel $\ker(\phi_p) = p\loc{p}$.



%%%%%%%%%%%%%%%%%%%%%%%%%%%%%%%%%%%%%%%%%%%%%%%%%%%%%%%%%%%%%%%%%%%%%%%%%%%%%%%%


%%%%%%%%%%%%%%%%%%%%%%%%%%%%%%%%%%%%%%%%%%%%%%%%%%%%%%%%%%%%%%%%%%%%%%%%%%%%%%%%



%\albertimportant{In case of polynomial rings $\cR[X]$, do we need to specify a degree bound? Technically not, but having a single size parameter could allow for wiggle room of degrees, making some descriptions and analysis a bit more tricky. Also Zip+ is designed having in mind that there's a fixed degree bound. We could have $B$ be a finite subset of $\cR$ instead of a bit bound}

%%%%%%%%%%%%%%%%%%%%%%%%%%%%%%%%%%%%%%%%%%%%%%%%%%%%%%%%%%%%%%%%%%%%%%%%%%%%%%%%
%%%%%%%%%%%%%%%%%%%%%%%%%%%%%%%%%%%%%%%%%%%%%%%%%%%%%%%%%%%%%%%%%%%%%%%%%%%%%%%%


%The following lemma on division by monic polynomials is proved in \cref{lem:monic_lift_Q_to_Z_appendix}.
%\begin{lemma}[Divisions of integer polynomials by monic integer polynomials]
%\label{lem:monic_lift_Q_to_Z}
%Let $f(X)\in \ZZ[X]$ and let $g(X)\in \ZZ[X]$ be monic. If $f(X)\in (g)_{\QQ[X]} \subseteq \QQ[X]$, i.e.\ there exists $h(X)\in \QQ[X]$ such that $f(X)=g(X)\cdot h(X)$, then in fact $f(X)\in (g)_{\ZZ[X]}\subseteq \ZZ[X]$ (equivalently, $h(X)\in \ZZ[X]$).
%\end{lemma}

%%%%%%%%%%%%%%%%%%%%%%%%%%%%%%%%%%%%%%%%%%%%%%%%%%%%%%%%%%%%%%%%%%%%%%%%%%%%%%%%
%%%%%%%%%%%%%%%%%%%%%%%%%%%%%%%%%%%%%%%%%%%%%%%%%%%%%%%%%%%%%%%%%%%%%%%%%%%%%%%%

%%%%%%%%%%%%%%%%%%%%%%%%%%%%%%%%%%%%%%%%%%%%%%%%%%%%%%%%%%%%%%%%%%%%%%%%%%%%%%%%
%%%%%%%%%%%%%%%%%%%%%%%%%%%%%%%%%%%%%%%%%%%%%%%%%%%%%%%%%%%%%%%%%%%%%%%%%%%%%%%%




\subparagraph{Constrained linear codes.}

A \emph{linear code} over a field  $\KK$ is a linear subspace $\mc{C}\subseteq \KK^n$.
The \emph{block length} is $n$, and the \emph{dimension} $k=\dim(\mc{C})$ is the dimension of the subspace.
A \emph{generator matrix} $M\in \KK^{n\times k}$ for $\mc{C}$ is a matrix whose columns form a basis for $\mc{C}$, so $\mc{C}=\{Mw : w\in \KK^k\}$.
\albert{Define distance first and then minimum distance and so on} The \emph{relative distance} $\delta\in(0,1]$ of $\mc{C}$ is the minimum, over all distinct $\cc,\cc'\in\mc{C}$, of the fraction of coordinates in which $\cc$ and $\cc'$ differ.
For a vector $\vv\in \KK^n$, the \emph{relative Hamming distance} from $\vv$ to the code is $\dist(\vv,\mc{C})=\min_{\cc\in\mc{C}}\dist(\vv,\cc)$, where $\dist(\vv,\cc)$ denotes the fraction of coordinates in which $\vv$ and $\cc$ differ.

Given a linear code $\mc{C}\subseteq \KK^n$ with generator matrix $M\in \KK^{n\times k}$, the \emph{$J$-wise interleaving} of $\mc{C}$ is the code
\[
  \mc{C}^J \;:=\; \{WM^T : W\in \KK^{J\times k}\} \;\subseteq\; \KK^{J\times n}.
\]
A codeword in $\mc{C}^J$ can be viewed as a tuple $(v_1,\ldots,v_J)$ of $J$ codewords $v_i\in\mc{C}$.

\subparagraph{Indexed relations and oracles.}
\label{sec:zinc-crypto-polynomial-commitments}


An \emph{indexed relation} $\REL$ is a set of index-input-witness triples $(\idx,\inp;\wit)$. The \emph{language} associated with $\REL$ is
\(
  \LANG(\REL)\ \defeq\ \bigl\{(\idx,\inp)\ \bigm|\ \exists\,\wit \text{ such that } (\idx,\inp;\wit)\in\REL\bigr\}.
\)



An \emph{oracle} to a string or vector $\pi$ of elements from a ring $R$ is a black-box that allows querying any position $j$ and obtaining $\pi_j$. We denote it by $\oracle{\pi}$. 
An \emph{oracle to a polynomial} $f(\XX)$ over $R$ is a black-box that allows querying any point $\rr$ in the domain of $f$ and obtaining $f(\rr)$. We denote it by $\oracle{f}$.

We provide standard definitions for IOPs, PIOPs, and their properties in \cref{sec:IOPs}. 

\begin{remark}[Oracle typing assumption]
\label{rem:zinc-crypto-oracle_typing}
Soundness and knowledge soundness for PIOPs are required to hold only for malicious provers $\prover^*$ that send \emph{correctly typed} oracles---i.e., oracles that actually correspond to polynomials with the prescribed number of variables, degrees, and coefficient domains. Informally speaking, enforcing that a (computationally bounded) prover sends such oracles is the responsibility of the polynomial commitment scheme (PCS) or other cryptographic mechanisms (such as IOPP) used to compile the PIOP into a concrete argument system.
\end{remark}






% \subparagraph{Constrained codes.}
% A \emph{constrained code} is a subset of a linear code defined by additional linear (or affine) constraints on the message space.
% Concretely, given a linear code $\mc{C}=\{Mw : w\in \KK^k\}$ and constraint vectors $u_1,\ldots,u_T\in \KK^k$ with target values $a_1,\ldots,a_T\in \KK$, the constrained code is
% \[
%   \mc{C}_{u_1\to a_1,\ldots,u_T\to a_T}
%   \;:=\;
%   \bigl\{Mw : w\in \KK^k,\ \langle u_t,w\rangle = a_t \text{ for all } t\in[T]\bigr\}.
% \]
% The constrained code inherits all the terminology corresponding to $\cC$, i.e.\ dimension, block length, rate, minimum distance, etc.  The notion extends to interleaved codes: given $\mc{C}^J$ and matrix-valued constraints, the so-called \emph{constrained interleaved code} $\mc{C}^J_{u_1\to a_1,\ldots,u_T\to a_T}$ requires that each column of the message matrix $W$ satisfies the corresponding constraint \albert{Is this how we define it? }.
% For precise definitions in the setting of codes over $\QQ$ (with optional modular reduction), see \cref{def:constrained_interleaved_code_relation}.
% \albert{TODO: Add citation for constrained codes if a standard reference exists.}


%\albert{Review this and make sure that this terminology and abuse of notation is okay and used correctly throughout the paper. Check if we need to add some remark about extractors extracting things in the subset code or if things are clear and non-ambiguous already. } Sometimes we will speak of codes $\cC$ over a subring $R$ of a field $\Frac{R}$, rather than over a field $\KK$. Concretely, $R$ for us will often be a ring of polynomials over a field (which is a subring of The corresponding field of rational functions). In this case, we abuse the notation and understand $\cC$ to be the subset of a code over the field $\Frac{R}$---i.e., $\cC \subseteq \Frac{R}^n$--- consisting of encodings whose message is a vector over $R$. 

\subparagraph{IOPs of proximity.}
An \emph{IOP of proximity (IOPP)} is an IOP variant designed to test whether an input oracle is \emph{close} to a member of some set (typically a code or a constrained code), rather than exactly a member.

\begin{definition}[Proximity relation]
\label{def:zinc-crypto-proximity_relation}
The \emph{proximity relation} $\REL_{\mathsf{prox}}$ for (possibly constrained) codes $\cC$ over $\KK$ and proximity parameters $\beta$ is defined as follows:
\[
  \REL_{\mathsf{prox}}
  \;:=\;
  \bigl\{(\idx,\inp;\wit)\ \bigm|\ \idx = (\mathcal{C},\beta),\ \inp = \oracle{\vv},\ \wit = \cc \in \mathcal{C},\ \dist(\vv,\cc) \le \beta\bigr\}.
\]
\end{definition}

\begin{definition}[Interactive Oracle Proof of Proximity (IOPP)]
\label{def:zinc-crypto-iopp}
Let $\idx=(\cC, \beta)$ be an index for $\REL_{\mathsf{prox}}$. 
An \emph{interactive oracle proof of proximity (IOPP)} for $\idx$ is an IOP for the proximity relation $\REL_{\mathsf{prox}}$ with  index restricted to $\idx = (\mathcal{C},\beta)$ (\cref{def:zinc-crypto-proximity_relation}), and with completeness only guaranteed when $\vv$ is in the code $\cC$, rather than $\beta$-close to it.
\begin{comment}
Concretely, it is a public-coin interactive protocol $(\prover,\verifier)$ in which:
\begin{itemize}
  \item $\prover$ and $\verifier$ receive $(\idx, \inp)$ as inputs, where  $\inp = \oracle{\vv}$ for some $\vv \in \KK^n$. The honest prover additionally receives a witness $\cc \in \mathcal{C}$ such that $\dist(\vv,\cc) \le \beta$.
  \item In each round, $\prover$ sends oracles to strings, and $\verifier$ sends a uniformly random challenge.
  \item At the end of the interaction, $\verifier$ queries the input oracle $\oracle{\vv}$ and the prover's oracle messages at positions of its choosing and outputs $\mathsf{accept}$ or $\mathsf{reject}$.
\end{itemize}
\noindent \textbf{Completeness.} For every $\vv \in \mathcal{C}$, there exists a prover $\prover$ such that
\[
  \Pr\bigl[\langle\prover(\idx,\inp,\wit),\verifier(\idx,\inp)\rangle = \mathsf{accept}\bigr]\ \ge\ 1-\varepsilon_{\mathsf{cmp}}.
\]
\noindent \textbf{Soundness.} For every $\vv \in \KK^n$ with $\dist(\vv, \mathcal{C}) > \beta$ and every prover $\prover^*$,
\[
  \Pr\bigl[\langle\prover^*,\verifier(\idx,\inp)\rangle = \mathsf{accept}\bigr]\ \le\ \varepsilon_{\mathsf{snd}}.
\]
\albert{Missing RBR Knowledge Soundness}
\end{comment}
\end{definition}



%%%%%%%%%%%%%%%%%%%%%%%%%%%%%%%%%%%%%%%%%%%%%%%%%%%%%%%%%%%%%%%%%%%%%%%%%%%%%%%%
%%%%%%%%%%%%%%%%%%%%%%%%%%%%%%%%%%%%%%%%%%%%%%%%%%%%%%%%%%%%%%%%%%%%%%%%%%%%%%%%

Extended definitions, including formal definitions of IOPs, PIOPs, PIORs, and Mutual Correlated Agreement, are given in \cref{s:extended_preliminaries}.
